\chapter{Introdução}
\label{chp:introducao}

\section{Contextualização}

A crescente demanda por processamento e análise de grandes volumes de dados
impulsionou a adoção de arquiteturas de \textit{Data Lake}, que permitem o
armazenamento centralizado de dados brutos em diferentes formatos:
estruturados, semiestruturados e não estruturados, sem necessidade de
transformação prévia \citep{tast2026comparative}. Nesse sentido, o armazenamento é
desacoplado do processamento, favorecendo flexibilidade e escalabilidade. Sendo
amplamente utilizado em pipelines de aprendizado de máquina e análise de dados
em larga escala, o \textit{Data Lake} atua como repositório central para dados de
diferentes fontes, como vídeos, imagens e documentos.

Em cenários de nuvem híbrida, organizações que buscam mitigar riscos de soberania
de dados e de dependência de fornecedor (\textit{vendor lock-in}) optam por manter
parte da infraestrutura de armazenamento \textit{on-premises}, replicando
funcionalidades oferecidas por provedores como Amazon Web Services (AWS), Microsoft
Azure e Google Cloud Platform (GCP). A camada de armazenamento de objetos é
componente central nessas arquiteturas, e o protocolo S3 da AWS consolidou-se como
padrão de fato para interoperabilidade entre sistemas \citep{tast2026comparative}.

Nesse contexto, o MinIO tornou-se uma solução amplamente utilizada para
armazenamento de objetos compatível com a API S3 em arquiteturas de \textit{Data Lake}
híbridas \citep{abdel2026pipeline}. Ao longo dos anos, a empresa responsável pelo
projeto promoveu mudanças progressivas em seu modelo de distribuição e licenciamento:
em 2021, a licença foi alterada de Apache 2.0 para GNU AGPLv3; em maio de 2025, a
interface de administração web foi removida da edição comunitária; em outubro de
2025, as distribuições binárias pré-compiladas foram descontinuadas, passando a
edição comunitária a ser distribuída apenas como código-fonte; e em fevereiro de
2026, o repositório público foi oficialmente marcado como não mantido
(\textit{THIS REPOSITORY IS NO LONGER MAINTAINED}), com correções de segurança
passando a ser avaliadas caso a caso, sem garantias \citep{minio2026repo,devpro2025,faun2026}.
O suporte ativo, os binários e as funcionalidades completas de administração foram
concentrados no produto comercial MinIO AIStor, disponível por assinatura.

Esse cenário expõe organizações que basearam sua camada de armazenamento de objetos
no MinIO a riscos concretos de segurança, de conformidade e de sustentabilidade
tecnológica, tornando relevante a avaliação de alternativas viáveis com
licenciamento aberto e compatibilidade com a API S3, além da governança de dados,
para mitigar riscos de dependência tecnológica e de \textit{vendor lock-in}.

\section{Proposta}

Este trabalho propõe um estudo experimental para avaliar comparativamente ao menos
três soluções de armazenamento de objetos compatíveis com a API S3, utilizando o
MinIO como \textit{baseline} de referência, sendo o RustFS uma das alternativas
obrigatoriamente incluídas. As demais soluções serão selecionadas durante o
levantamento bibliográfico, priorizando projetos com licenciamento aberto e adoção
ativa pela comunidade.

O objetivo geral é avaliar comparativamente alternativas ao MinIO para armazenamento
de objetos em arquiteturas de \textit{Data Lake} híbridas, considerando aspectos de
desempenho, compatibilidade com a API S3, interoperabilidade com provedores de
nuvem, custos, licenciamento, sustentabilidade e riscos tecnológicos.

São objetivos específicos deste trabalho:

\begin{itemize}
    \item Identificar soluções de armazenamento de objetos compatíveis com a API S3
          que possam constituir alternativas ao MinIO em arquiteturas híbridas;
    \item Caracterizar os requisitos técnicos relevantes para armazenamento de
          objetos em arquiteturas de \textit{Data Lake} híbridas;
    \item Implementar um ambiente experimental padronizado para implantação e
          avaliação das soluções selecionadas;
    \item Investigar a compatibilidade das soluções com a API S3 e sua
          interoperabilidade com serviços de nuvem (AWS, Azure e GCP);
    \item Avaliar o desempenho e o consumo de recursos das soluções em cenários
          controlados com conjuntos de dados de diferentes características;
    \item Analisar comparativamente as soluções segundo custos, modelo de
          licenciamento, maturidade do projeto e riscos tecnológicos.
\end{itemize}

\section{Justificativa}

As mudanças progressivas no modelo de distribuição e licenciamento do MinIO
representam um risco concreto para organizações públicas e privadas que adotaram
a solução como componente central de sua infraestrutura de dados. A descontinuação
das distribuições binárias e o congelamento do desenvolvimento comunitário em 2025,
seguidos pelo arquivamento do repositório público em fevereiro de 2026
\citep{faun2026,minio2026repo}, comprometem a previsibilidade operacional de
ambientes que dependem da solução, especialmente em contextos de nuvem híbrida onde
soberania de dados e controle sobre a infraestrutura são requisitos críticos.

A ausência de estudos comparativos publicados que avaliem alternativas ao MinIO
nesse cenário específico --- considerando soluções emergentes como o RustFS e
aspectos como compatibilidade com provedores de nuvem, sustentabilidade do
licenciamento e riscos de dependência tecnológica --- justifica a relevância
acadêmica desta pesquisa \citep{tast2026comparative}. Os resultados poderão
orientar equipes de engenharia de dados na seleção de tecnologias mais sustentáveis,
contribuindo para a redução de riscos de segurança e de \textit{vendor lock-in}.
